\section{Conclusion et perspectives}
\label{sec:conclusion}

\begin{resume}
La recherche multic\oe{}ur par vagues est livr\'ee, test\'ee, mesur\'ee et
active dans le bot. Elle apporte $+8$ \`a $+26\,\%$ de d\'ebit selon la
position, sans d\'egradation de la queue de latence ni de la qualit\'e de
recherche, et sans r\'egression du self-play. Le plafond observ\'e est
intrins\`eque : l'inf\'erence r\'eseau est le goulot. La suite naturelle n'est
donc pas davantage de workers, mais l'assainissement de l'appel \`a
l'\'evaluateur et, seulement apr\`es, la production continue.
\end{resume}

\subsection{Bilan}

Le chantier a transform\'e une boucle strictement s\'erielle en une
architecture \`a deux phases, avec une surface concurrente volontairement
r\'eduite : quatre \'etats de n\oe{}ud atomiques, un compteur de virtual loss,
des snapshots de table sous verrous ray\'es, et une garde RAII qui poss\`ede
le chemin r\'eserv\'e. La d\'emonstration de correction repose sur des
invariants v\'erifi\'es apr\`es chaque recherche, des entrelacements forc\'es
et des tests de mutation, faute de ThreadSanitizer sous MSVC.

Sur le plan des performances, le verdict est mesur\'e et le m\'ecanisme est
compris : deux workers capturent l'essentiel du gain, huit l'emportent
l\'eg\`erement sur les m\'edianes et sur la queue, et la zone gagnante est la
finale. Sur le plan de la qualit\'e, la non-inf\'eriorit\'e est acquise sur
2\,500 puzzles, avec une r\'eserve assum\'ee : la r\'ep\'etition du candidat
n'a pas \'et\'e faite.

\subsection{Ce que la production continue peut reprendre}

L'\'etape~2 de la conception, la production continue, ferait descendre les
workers pendant les inf\'erences et les remont\'ees, au lieu de les arr\^eter
\`a une barri\`ere. Tous ses composants existent d\'ej\`a : contextes priv\'es,
r\'esultats immuables, garde RAII, propri\'et\'e \code{Pending}, cache
thread-safe, \'evaluateur contr\^olable, ex\'ecuteur persistant. Ce qui
manque est la synchronisation des statistiques et de la structure d'arbre
pendant que d'autres threads lisent.

Les mesures de ce rapport permettent d'en estimer l'enjeu : l'attente de
barri\`ere ne repr\'esente que 1,5 \`a 2,3\,\% du temps, et le remplissage
plafonne d\'ej\`a vers 7. Autrement dit, la production continue ne peut pas
cr\'e\'er beaucoup plus de parall\'elisme utile \`a lot de 8 ; son int\'er\^et
d\'epend de la possibilit\'e de former des lots plus gros, ce qui est pr\'ecis\'ement
la question du chantier suivant.

\subsection{Assainir l'appel \`a l'\'evaluateur}

L'\'etude \code{2026-09-19-cout-calcul-cpu-gpu.md} a mis en \'evidence un
\'ecart inattendu : le banc self-play \'evalue huit positions en 3,0\,ms,
alors que la recherche paie 10 \`a 14\,ms par appel \`a remplissage
comparable. La diff\'erence est attribu\'ee \`a la forme variable des lots et
aux allocations par appel c\^ot\'e ONNX Runtime. Le test T1b de cette \'etude
(forme de lot fixe, buffers r\'eutilis\'es) est l'exp\'erience la plus
rentable identifi\'ee \`a ce jour, avec un potentiel de facteur 3 \`a 5 sur
l'appel. Si elle se confirme, la production continue devient beaucoup plus
int\'eressante, parce que de plus gros lots deviendront utiles. C'est
l'ordre logique : d'abord l'appel, ensuite l'ordonnancement.

\subsection{Autres pistes consign\'ees}

Mesurer le cas chaud (arbre r\'eutilis\'e, table de 4\,000\,000 d'entr\'ees)
pour chiffrer la part de simulations qui \'evitent le r\'eseau en partie
r\'eelle. Revoir la taille de table maintenant que la recherche tourne \`a
des budgets plus \'elev\'es. Envisager FP16 ou TensorRT pour le r\'eseau,
apr\`es T1b. Pour le self-play, une file d'inf\'erence centrale pourrait
lisser les vagues GPU, mais les mesures montrent qu'il est d\'ej\`a efficace
par parties ind\'ependantes.

\subsection{Mot de la fin}

Ce chantier aura surtout montr\'e une chose : la performance d'un moteur
d'\'echecs neuronal ne se joue pas l\`a o\`u l'intuition la place. Le
multic\oe{}ur \'etait la suite logique du batching, et il rend bien ce qu'il
peut rendre, environ 10\,\% de m\'ediane, parce que 97\,\% du temps est
d\'ej\`a dans le GPU. La vraie r\'eserve de performance est dans la
plomberie de l'appel \`a l'\'evaluateur, pas dans le nombre de c\oe{}urs.
Mesurer avant d'optimiser, et mesurer le r\'egime r\'eel plut\^ot que le
r\'egime commode, aura \'et\'e la boussole de tout le chantier.
